\section{Method}
\label{sec:method}

\subsection{Five-Stage Workflow}

The refactored paper module adopts a five-stage workflow adapted from the design pattern used in ARA: idea formation, experiment derivation, writing, revision, and build export. In this repository, those stages map to concrete directories and artifacts rather than abstract prompts. The idea stage defines research questions and scope constraints. The experiment stage generates structured CSV summaries and figures from the normalized incident corpus and from protocol-study constants already supported by repository reports. The writing stage contains a blueprint, bibliography, modular \LaTeX{} sections, and figure assets. The revision stage records weaknesses in the previous draft and the changes introduced in this refactor. Figure~\ref{fig:framework} shows the corresponding research workflow that connects incident collection to paper synthesis.

\begin{figure}[t]
\centering
\includegraphics[width=0.98\textwidth]{research_framework.png}
\caption{Refactored paper workflow. The pipeline begins with a broad incident corpus, narrows through normalization and protocol screening, and only then promotes selected cases into replay engineering and paper claims.}
\label{fig:framework}
\end{figure}

\subsection{Evidence Layers}

The workflow is built around three evidence layers. The first is a normalized incident dataset loaded from `data/processed/incidents\_normalized\_latest.csv`, which currently contains 2{,}036 rows and 2{,}035 unique incident identifiers, of which 1{,}378 are labeled as DeFi. The second is a protocol evidence layer composed of five deep audits and a 24-protocol pure on-chain verification study. The third is a replay layer consisting of twelve shortlisted incidents, with concrete replay logic already encoded for Euler Finance, BonqDAO and AllianceBlock, and Lendf.Me.

We intentionally distinguish between \emph{confirmed exploitability}, \emph{conditional exploitability}, and \emph{replay engineering progress}. Confirmed exploitability means that the repository contains high-confidence evidence of a pure on-chain attack path under the study's threat model. Conditional exploitability means that the issue depends on downstream integrations, special market states, or missing runtime confirmation. Replay engineering progress is weaker still: it indicates whether the exploit path has been encoded as executable logic in tests, regardless of whether archive RPC validation has been completed.

\begin{table}[t]
\centering
\caption{Evidence layers and the kinds of claims they support.}
\label{tab:evidence-layers}
\begin{tabular}{p{2.2cm}p{4.0cm}p{5.3cm}}
\toprule
\textbf{Layer} & \textbf{Primary artifacts} & \textbf{Supported claims} \\
\midrule
Incident corpus & Normalized incident CSV, yearly counts, family taxonomy, loss summaries & Broad ecosystem trends, dominant categories, shortlist construction, evidence that public feeds are heterogeneous \\
Protocol evidence & Deep audits, triaged findings, 24-protocol verification report & High-confidence statements about pure on-chain exploitability under a specific threat model \\
Replay evidence & Incident cards, replay tests, fork anchors, encoded exploit logic & Mechanism-level claims about exploit preconditions and reproduction feasibility \\
\bottomrule
\end{tabular}
\end{table}

Table~\ref{tab:evidence-layers} summarizes the paper's central discipline: every claim must be attached to the layer that actually supports it. This prevents a common failure mode in security writing, where scanner findings are discussed as if they were reproduced exploits or where incident counts are discussed as if they were code-level proof.

\subsection{Scoring and Derived Metrics}

Replay work is prioritized using a pragmatic scoring rule rather than a claim of theoretical severity. For an incident $i$, we define a replayability prioritization score as
\[
S(i) = \alpha_f F(i) + \alpha_l \log(1 + L(i)) + \alpha_r R(i) + \alpha_c C(i),
\]
where $F(i)$ weights historically replayable families such as oracle manipulation, flash-loan-assisted logic abuse, and reentrancy; $L(i)$ is the reported loss in U.S. dollars; $R(i)$ indicates whether a public exploit path or post-mortem exists; and $C(i)$ represents chain and contract identifiability. This score is used operationally for shortlist construction, not as a normative security metric.

For the paper, two aggregate metrics are more important. The first is replay conversion,
\[
\mathrm{ReplayConversion} = \frac{N_{\mathrm{replay\ logic}}}{N_{\mathrm{shortlisted}}},
\]
which measures how many shortlisted incidents have already been turned into concrete replay logic. The second is protocol exposure ratio,
\[
\mathrm{ExposureRatio} = \frac{N_{\mathrm{conditional}} + N_{\mathrm{confirmed}}}{N_{\mathrm{verified}}},
\]
which measures how much of the verified protocol set remains exposed to either direct or integration-dependent pure on-chain attack paths. In the current artifact set, these values are $3/12 = 0.25$ and $4/24 \approx 0.167$, respectively.

\subsection{Threat Model and Confidence Boundaries}

The protocol-study layer uses a deliberately narrow threat model. An attack is treated as pure on-chain only if it can be executed in a single transaction or a short block window without relying on external compromise, extraordinary governance capture, or unrealistic market conditions. This choice biases the study toward precision rather than recall. It undercounts some real incidents, but it avoids conflating every observed loss with directly replayable contract exploitability. The replay layer inherits the same discipline: a scaffold or encoded exploit path is not described as fully verified unless archive RPC execution has actually succeeded.
